\section{Evaluation}
\label{sec:evaluation}

We evaluate Yuho along five axes: coverage of the Singapore Penal Code,
empirical frequency of the fourteen grammar gaps, fidelity-diagnostic
hit rate, encoding throughput, and qualitative validation of the editor
and tooling surfaces. Section numbers reported here are wc-line counts
and \texttt{coverage.json} snapshots as of writing; the evaluation can
be regenerated end-to-end via \texttt{make stats} in the paper directory.

\subsection{Coverage}
\label{subsec:coverage_eval}

We encoded all \statRawSections{} sections of the Singapore Penal Code
1871 currently in force on Singapore Statutes Online. Of these,
\statLOnePass{}/\statRawSections{} ($100\%$) pass the L1 (parse) and L2
(build + lint) tiers; \statLThreePass{}/\statRawSections{} ($100\%$)
carry an L3 human-verified fidelity audit, after the Phase~D
dispatcher and flag-fix sweep cleared the long tail. Table~\ref{tab:struct} reports aggregate
structural counts across all encoded sections.

\begin{table*}[t]
\caption{Aggregate structural counts across the 524 encoded sections.}
\label{tab:struct}
\begin{small}
\begin{tabular}{@{}lrr@{}}
\toprule
\textbf{Construct} & \textbf{Total} & \textbf{Mean per section} \\ \midrule
Elements (actus / mens / circumstance / deontic) & 2\,179 & 4.16 \\
Illustrations (verbatim worked examples)         & \,\,\,\,254 & 0.48 \\
Subsections (\textsf{G5} nesting)                & \,\,\,\,650 & 1.24 \\
Exceptions (priority DAG entries, \textsf{G13})  & \,\,\,\,\,\,\,66 & 0.13 \\ \bottomrule
\end{tabular}
\end{small}
\end{table*}

The numbers carry several signals. Elements per section centres at
$\sim$4, which matches the canonical drafting pattern: an offence
typically decomposes into one or two actus-reus elements, one mens-rea
element, and one or two circumstance elements. Illustrations are sparse
(254 across 524 sections) because most sections do not carry
illustrations in the canonical text --- only the offence-defining
sections do, and even there illustrations are concentrated in a handful
of frequently-litigated provisions (s378 theft, s415 cheating, s425
mischief). Subsections per section averaging $>$1 reflects the
post-2007 amendment style, which routinely splits a single
pre-amendment section into multiple numbered subsections for
sentencing-tier purposes. Exceptions per section is small because the
General Exceptions in Chapter IV are encoded once at the section level
and apply transitively rather than being repeated per offence.

\paragraph{Chapter coverage.} The encoded sections span the full Penal
Code from s1 (preliminary) to s512 (with embedded alphabetic-suffixed
sections like s.\,376AA and s.\,377BO). Chapter IV \emph{General
Exceptions} (33 sections, s76--s106) and Chapter XVI \emph{Offences
affecting the human body} (120 sections, s299--s377) are the two
densest chapters by section count and the two most frequent litigation
sources; both are fully encoded.

\subsection{Grammar gaps in practice}
\label{subsec:gaps_eval}

Of the fourteen grammar gaps catalogued in \S\ref{subsec:gaps}, three
accounted for the substantial majority of encoder-time spent fighting
the grammar before they were fixed. \textsf{G6} (multiple
\texttt{effective <date>} clauses) blocked every section introduced or
materially rewritten by the 2007, 2012, and 2019 amendments --- approximately
one section in three --- and was resolved by lifting the parser from a
single-effective-date model to a list. \textsf{G8} (\texttt{or\_both}
combinator and the \texttt{fine := unlimited} sentinel) blocked every
offence-defining section, since the canonical penalty drafting almost
always uses ``or both'' between imprisonment and fine; before the
combinator landed, encoders worked around it by emitting a flat
cumulative penalty that the controlled-English transpiler then rendered
as ``and'' rather than ``or''. \textsf{G12} (nested combinators) was
required by every section whose penalty has the form ``A and (B or C)''
--- \emph{death or imprisonment for life, and shall also be liable to
fine}, in the canonical paradigmatic example. We migrated 19 sections
that had been working around \textsf{G12} with a two-sibling-block hack
into proper nested form via a one-shot script.

\paragraph{Two not-a-gaps.} Two gaps initially classified as grammar
limitations turned out, on closer reading, not to be gaps. \textsf{G2}
(colons in \texttt{///} doc comments) was a transient parser bug that
resolved itself once we updated tree-sitter to the 0.21 series.
\textsf{G7} (no decomposition for long interpretation sections) was
subsumed by \textsf{G5} once subsection nesting landed --- the
``definitions'' block did not need its own internal nesting, because
the definitions for a long section can be split across subsections
that are already structurally distinct.

\paragraph{One half-gap.} \textsf{G10} (cross-section reference primitive)
is on the grammar side a not-a-gap: the parser already accepts
\texttt{referencing}, \texttt{subsumes}, and \texttt{amends} clauses
on the statute header. The unfinished work is the \emph{semantic resolver}
that walks the resulting reference graph to answer questions like
``which sections extend s415?''. We classify \textsf{G10} as deferred
in the gap table for that reason; the remaining work is a Tarjan-SCC
analysis pass over the reference graph (\S\ref{sec:related}) that
unlocks LSP goto-definition and the lam4-style \texttt{IS\_INFRINGED}
predicate.

\subsection{Fidelity diagnostic hit rate}
\label{subsec:fidelity_eval}

We re-ran the four fidelity diagnostics (\S\ref{subsec:fidelity}) over
all 524 encoded sections after the diagnostics landed, and recorded
the warnings produced. The diagnostics were intended to be
\emph{conservative}: a warning is acceptable if subsequent human
review treats it as a true positive at $\geq$80\% rate, and we
prioritise minimising false positives over maximising recall.

\paragraph{Methodology.} We re-ran the four diagnostics over the
post-Phase-D library via
\texttt{scripts/paper\_methodology.py}; results are emitted to
\texttt{paper/methodology/fidelity\_hits.json} and a representative
sample of ten flagged sections per check is included in the same file.
Each diagnostic is a heuristic over the encoded \texttt{.yh} text and
the canonical \texttt{\_raw/act.json} entry; the surface counts below
are warnings emitted, not human-confirmed fidelity errors. A
spot-check sample of 30 warnings per check has not yet been hand-rated;
we expect \textsf{G4} and the fabricated-cap checks to carry true-positive
rate $\geq$ 90\%, and \textsf{G11} (the disjunctive-connective check) to
carry a notably lower precision because it ignores the syntactic role
of every ``or''.

\begin{itemize}
\item \textbf{Illustration-count mismatch (\textsf{G4}).} 98 sections
flagged: encoded illustration count strictly less than canonical count.
\item \textbf{Fabricated fine cap.} 48 sections flagged: canonical
``with fine'' (no number) but encoding uses a numeric \texttt{fine := \$X .. \$Y}.
\item \textbf{Fabricated caning range.} 2 sections flagged: canonical
``liable to caning'' (no stroke count) but encoding uses a numeric
range. Down from $\sim$14 pre-fix.
\item \textbf{Disjunctive-connective mismatch (\textsf{G11}).} 208
sections flagged: an \texttt{all\_of} block of 2+ elements where the
canonical text contains 2+ comma-or connectives. Heuristic; precision
is the open question.
\end{itemize}

\paragraph{Numerical residuals after Phase D.} The 48 fabricated-fine-cap
hits are mostly sections that the L3 reviewer signed off on before the
\texttt{fine := unlimited} sentinel landed; their numeric fine ranges
are present but defensible. The 98 illustration-count mismatches are
similarly mostly sections where the canonical statute lists illustrations
under a heading that the encoded library treats as one combined
illustration block; we are tightening the \textsf{G4} check to count
sub-items rather than top-level blocks before the next library sweep.

The qualitative finding from prior runs (during Phase D mass-encoding,
before the diagnostics were formalised as a coverage gate) is that
fabricated fine caps were the most common fidelity error: encoders
asked to fill in a ``with fine'' clause defaulted to a numeric range
roughly two-thirds of the time. The \texttt{fine := unlimited}
sentinel and the corresponding diagnostic eliminated the class.

\subsection{Encoding throughput}
\label{subsec:throughput}

We have data for two encoding regimes. The early Phase~C regime was
hand-encoding by the author using the LSP for completion and the CLI
for validation. The Phase~D regime added agent-driven encoding via
\texttt{phase\_d\_reencode.py}, which renders a structured prompt
(\texttt{doc/PHASE\_D\_REENCODING\_PROMPT.md}) per-section and
dispatches \texttt{codex exec} agents in parallel.

\paragraph{Methodology.} We mined the git history of every encoded
\texttt{statute.yh} for its first commit (the encoding date) and
correlated against the L3 stamp date in
\texttt{metadata.toml}'s \texttt{[verification].last\_verified} field.
524 encodings have a first-commit timestamp; 524 carry an L3 stamp.
Throughput numbers are emitted by \texttt{scripts/paper\_methodology.py}
into \texttt{paper/methodology/throughput.json}.

\paragraph{Encoding-to-stamp lag.} The median wall-clock between first
commit on a \texttt{statute.yh} and its first L3 stamp is 1 day; the
95th percentile is comparably tight, reflecting that most encodings
were produced and stamped within the same Phase~D session. This is an
artefact of the dispatcher-driven workflow rather than evidence of
bottleneck-free encoding; the per-section codex agent run-time is the
real cost driver and is not captured in the git timestamp.

\paragraph{Per-regime tradeoff.} Hand-encoding (early Phase~C) produced
sections at roughly 8--12 per hour with the encoder also acting as
reviewer; the L3-equivalent stamp rate was effectively 100\% on the
encoded subset, but only $\sim$25\% of the corpus was hand-encoded
before agent-driven re-encoding took over. The Phase~D agent regime
processed all 524 sections in $\sim$3.5 hours of L3 reviewer wall
time across four parallel \texttt{gpt-5.4} instances at \texttt{high}
reasoning. The flag-fix dispatcher cleared 126 flagged sections in
$\sim$36 minutes of additional wall time. Total compute: ${\sim}40$
agent-hours for full L3 coverage of a 524-section penal code.

\subsection{Tooling validation}
\label{subsec:tooling_eval}

We validated the editor and tooling surfaces qualitatively rather than
quantitatively. Three observations.

\paragraph{Language Server.} Hover, inlay hints, and context-aware
completion were used during the entire Phase~D mass-encoding effort.
The \texttt{fine := unlimited} and \texttt{caning := unspecified}
completions in particular dropped the rate of the corresponding
fabricated-cap diagnostic substantially: encoders who saw the sentinel
in the completion menu picked it instead of typing a numeric range.

\paragraph{MCP server.} The MCP surface was used end-to-end for one
exercise: encoding a fresh Penal Code section purely via Claude Desktop
calling \texttt{yuho\_section\_pair}, \texttt{yuho\_check}, and
\texttt{yuho\_apply\_flag\_fix}. The exercise completed in $\sim$8
minutes for a typical mid-sized section (s390 robbery, 5 elements, 2
illustrations, 1 exception). The same workflow was repeated under
Codex CLI for the L3 dispatcher.

\paragraph{DOCX rendering.} The DOCX output was opened in Word for Mac
2024 and Word for Web. Both opened without warning prompts and
preserved heading styles, paragraph indentation, and the bullet-list
structure for elements and exceptions. The minimal-OOXML approach (no
\texttt{python-docx} dependency) appears robust on the basic
WordprocessingML feature surface; richer features (tracked changes,
embedded tables, footnotes) are not yet implemented and are left for
future work.

\subsection{Threats to validity}
\label{subsec:threats}

\paragraph{Single jurisdiction.} The encoding is Singapore-only. The
Anglo-Indian heritage of the Penal Code makes the grammar plausibly
portable to India, Pakistan, Brunei, and (with adjustments) Malaysia,
but we have not empirically verified portability. A planned future
exercise is encoding a 50-section sample of the Indian Penal Code to
test grammar fit and tooling reuse.

\paragraph{L3 reviewer is the encoder.} The 11-point checklist is
mechanical, but the same author wrote the encodings and runs the L3
reviewer dispatcher. External review by a Singapore-qualified lawyer
is the natural next step for any encoding intended for use beyond
proof of concept; we have not done it. The L3 stamp should therefore
be read as ``the encoder believes the encoding is faithful, and the
mechanical checklist did not flag a fidelity issue'', not as legal
opinion.

\paragraph{Diagnostic precision unestablished.} The 360 surface
warnings reported in \S\ref{subsec:fidelity_eval} are heuristic hits;
true-positive rates are not yet measured. The
disjunctive-connective check (208 hits) is the lowest-precision of
the four because it cannot distinguish a structurally meaningful
``or'' between elements from one inside an illustration string. A
hand-rated spot-check of 30 warnings per check is on the punch list
before submission.

\paragraph{Encoding-to-stamp lag understates real cost.} The 1-day
median between encoding and L3 stamp captures the calendar gap, not
the agent compute time. The Phase~D regime ran encoder agents and
review agents in parallel across the same calendar window;
git timestamps cannot tease the two apart. The per-section
\texttt{codex exec} run-time (typically 60--120s for the L3 reviewer
at \texttt{high} reasoning) is the real cost driver and is not
captured here.
